\chapter{Annexes Techniques et Bibliographie}
\label{chap:annexes}

\section{Spécification des Endpoints de l'API RESTful}

Le tableau~\ref{tab:api_endpoints_annex} documente la spécification formelle des principaux points d'accès exposés par le serveur Flask de Perceptra. Toutes les requêtes et réponses échangent des données sérialisées au format JSON, à l'exception des routes de téléversement et de téléchargement binaire.

\begin{table}[H]
\centering
\rowcolors{2}{LightGray}{white}
\scriptsize
\begin{tabularx}{\textwidth}{>{\bfseries}p{1.4cm}p{3.5cm}p{2.1cm}X}
\toprule
\rowcolor{TableHeader}
\color{white}Méthode & \color{white}Point d'Accès & \color{white}Auth. Requise & \color{white}Description \& Paramètres Clés \\
\midrule
\texttt{POST} & \texttt{/api/auth/register} & Non & Inscription d'un utilisateur (\texttt{username}, \texttt{email}, \texttt{password}). \\
\texttt{POST} & \texttt{/api/auth/login} & Non & Authentification et émission de la paire de jetons JWT. \\
\texttt{GET} & \texttt{/api/auth/profile} & JWT Bearer & Consultation des informations du compte utilisateur connecté. \\
\midrule
\texttt{GET} & \texttt{/api/images} & Optionnelle & Liste paginée des images (\texttt{page}, \texttt{limit}, \texttt{category\_id}, \texttt{favorites\_only}). \\
\texttt{POST} & \texttt{/api/images/upload} & JWT Bearer & Téléversement d'image (\texttt{multipart/form-data}, calcul SHA-256 et descripteurs). \\
\texttt{GET} & \texttt{/api/images/<id>/file} & Non & Téléchargement du flux binaire original. \\
\texttt{GET} & \texttt{/api/images/<id>/thumbnail}& Non & Téléchargement de la vignette allégée ($256 \times 256$ pixels). \\
\texttt{DELETE}& \texttt{/api/images/<id>} & JWT (Owner) & Suppression d'une image, de ses variantes et de ses descripteurs. \\
\texttt{POST} & \texttt{/api/images/batch-delete} & JWT Bearer & Suppression par lot d'une liste d'identifiants d'images. \\
\texttt{POST} & \texttt{/api/images/<id>/transform}& JWT Bearer & Retouches non destructives (\texttt{crop}, \texttt{resize}, \texttt{rotate}, \texttt{affine}). \\
\midrule
\texttt{GET} & \texttt{/api/categories} & Optionnelle & Arborescence des catégories sérialisée sous forme d'arbre récursif. \\
\texttt{POST} & \texttt{/api/categories} & JWT Bearer & Création d'une catégorie (\texttt{name}, \texttt{parent\_id}). \\
\midrule
\texttt{GET} & \texttt{/api/descriptors/<id>} & Optionnelle & Consultation des vecteurs de signatures numériques. \\
\texttt{POST} & \texttt{/api/descriptors/extract}& JWT Bearer & Déclenchement explicite de l'extraction sur une image donnée. \\
\midrule
\texttt{POST} & \texttt{/api/query/simple} & JWT Bearer & Recherche initiale par image (\texttt{image\_id}, \texttt{weights}, \texttt{metric}, \texttt{k}). \\
\texttt{POST} & \texttt{/api/query/feedback} & JWT Bearer & Soumission du feedback FRE (\texttt{query\_id}, \texttt{relevant\_ids}, \texttt{irrelevant\_ids}). \\
\midrule
\texttt{GET} & \texttt{/api/analytics/overview} & JWT Admin & Métriques globales du système (images, taille disque, requêtes, latence). \\
\bottomrule
\end{tabularx}
\caption{Spécification formelle des endpoints de l'API REST de Perceptra}
\label{tab:api_endpoints_annex}
\end{table}

\section{Schéma SQL DDL et Guide de Démarrage}

\begin{lstlisting}[language=SQL, basicstyle=\scriptsize\ttfamily, caption={Définition DDL des tables principales de Perceptra}]
CREATE TABLE users (
    id VARCHAR(36) PRIMARY KEY, username VARCHAR(50) NOT NULL UNIQUE,
    email VARCHAR(120) NOT NULL UNIQUE, password_hash VARCHAR(255) NOT NULL,
    role VARCHAR(20) NOT NULL DEFAULT 'user', created_at TIMESTAMP NOT NULL
);
CREATE TABLE categories (
    id VARCHAR(36) PRIMARY KEY, user_id VARCHAR(36) REFERENCES users(id) ON DELETE SET NULL,
    parent_id VARCHAR(36) REFERENCES categories(id) ON DELETE CASCADE,
    name VARCHAR(100) NOT NULL, path VARCHAR(500) NOT NULL, created_at TIMESTAMP NOT NULL
);
CREATE TABLE images (
    id VARCHAR(36) PRIMARY KEY, owner_id VARCHAR(36) REFERENCES users(id) ON DELETE SET NULL,
    category_id VARCHAR(36) REFERENCES categories(id) ON DELETE SET NULL,
    filename VARCHAR(255) NOT NULL, storage_key VARCHAR(500) NOT NULL UNIQUE,
    source_type VARCHAR(20) NOT NULL, processing_status VARCHAR(20) NOT NULL,
    mime_type VARCHAR(100) NOT NULL, file_size INT NOT NULL, width INT NOT NULL, height INT NOT NULL,
    checksum VARCHAR(64) NOT NULL, transformation_params JSON, created_at TIMESTAMP NOT NULL
);
CREATE INDEX ix_images_checksum ON images(checksum);
CREATE TABLE image_features (
    id VARCHAR(36) PRIMARY KEY, image_id VARCHAR(36) NOT NULL REFERENCES images(id) ON DELETE CASCADE,
    feature_type VARCHAR(50) NOT NULL, extractor_version VARCHAR(50) NOT NULL,
    feature_values JSON NOT NULL, parameters JSON, created_at TIMESTAMP NOT NULL,
    CONSTRAINT uq_image_feature_type_version UNIQUE (image_id, feature_type, extractor_version)
);
CREATE TABLE retrieval_queries (
    id VARCHAR(36) PRIMARY KEY, user_id VARCHAR(36) NOT NULL REFERENCES users(id) ON DELETE CASCADE,
    query_image_id VARCHAR(36) REFERENCES images(id) ON DELETE SET NULL,
    search_scope VARCHAR(30) NOT NULL, configuration JSON, status VARCHAR(20) NOT NULL, created_at TIMESTAMP NOT NULL
);
CREATE TABLE retrieval_iterations (
    id VARCHAR(36) PRIMARY KEY, query_id VARCHAR(36) NOT NULL REFERENCES retrieval_queries(id) ON DELETE CASCADE,
    iteration_number INT NOT NULL, retrieval_method VARCHAR(30) NOT NULL,
    feature_weights JSON, updated_query_vector JSON, feedback_count INT NOT NULL,
    execution_time_ms INT NOT NULL, results_count INT NOT NULL, created_at TIMESTAMP NOT NULL,
    CONSTRAINT uq_retrieval_iteration_query_number UNIQUE (query_id, iteration_number)
);
CREATE TABLE retrieval_results (
    id VARCHAR(36) PRIMARY KEY, iteration_id VARCHAR(36) NOT NULL REFERENCES retrieval_iterations(id) ON DELETE CASCADE,
    image_id VARCHAR(36) NOT NULL REFERENCES images(id) ON DELETE CASCADE,
    rank INT NOT NULL, similarity_score FLOAT NOT NULL, feature_scores JSON, created_at TIMESTAMP NOT NULL
);
CREATE TABLE relevance_feedbacks (
    id VARCHAR(36) PRIMARY KEY, iteration_id VARCHAR(36) NOT NULL REFERENCES retrieval_iterations(id) ON DELETE CASCADE,
    image_id VARCHAR(36) NOT NULL REFERENCES images(id) ON DELETE CASCADE,
    feedback_type VARCHAR(20) NOT NULL, created_at TIMESTAMP NOT NULL
);
\end{lstlisting}

\begin{lstlisting}[language=bash, basicstyle=\scriptsize\ttfamily, caption={Commandes d'initialisation et démarrage du projet Perceptra}]
# Cloner et initialiser le projet
git clone https://github.com/regisx001/Perceptra.git && cd Perceptra && cp .env.example .env
uv sync && uv run perceptra db upgrade
uv run perceptra dev     # Lancement backend Flask (port 5000)
cd web && bun install && bun run prepare && bun run dev  # Lancement frontend (port 5173)
# Alternative conteneurisee tout-en-un
docker compose up --build -d
\end{lstlisting}

\vspace{0.4cm}
\section{Bibliographie et Références Scientifiques}

\begin{thebibliography}{99}
\small
\bibitem{flickner1995qbic}
Flickner, M., et al. (1995). Query by image and video content: The QBIC system. \textit{IEEE Computer}, 28(9), 23--32.

\bibitem{pentland1996photobook}
Pentland, A., Picard, R. W., \& Sclaroff, S. (1996). Photobook: Content-based manipulation of image databases. \textit{IJCV}, 18(3), 233--254.

\bibitem{smith1996visualseek}
Smith, J. R., \& Chang, S. F. (1996). VisualSEEk: a fully automated content-based image query system. \textit{Proc. ACM Multimedia}, 87--98.

\bibitem{smeulders2000content}
Smeulders, A. W., et al. (2000). Content-based image retrieval at the end of the early years. \textit{IEEE TPAMI}, 22(12), 1349--1380.

\bibitem{stricker1995similarity}
Stricker, M. A., \& Orengo, M. (1995). Similarity of color images. \textit{SPIE}, Vol. 2420, 381--392.

\bibitem{haralick1973textural}
Haralick, R. M., et al. (1973). Textural features for image classification. \textit{IEEE Trans. SMC}, 3(6), 610--621.

\bibitem{ojala2002multiresolution}
Ojala, T., et al. (2002). Multiresolution gray-scale and rotation invariant texture classification with LBP. \textit{IEEE TPAMI}, 24(7), 971--987.

\bibitem{rocchio1971relevance}
Rocchio, J. J. (1971). Relevance feedback in information retrieval. \textit{The SMART Retrieval System}, Prentice-Hall, 313--323.

\bibitem{rui1998relevance}
Rui, Y., et al. (1998). Relevance feedback: a power tool for interactive content-based image retrieval. \textit{IEEE TCSVT}, 8(5), 644--655.

\bibitem{vaswani2017attention}
Vaswani, A., et al. (2017). Attention is all you need. \textit{NeurIPS}, 30, 5998--6008.

\bibitem{radford2021learning}
Radford, A., et al. (2021). Learning transferable visual models from natural language supervision (CLIP). \textit{ICML}, 8748--8763.

\bibitem{oquab2023dinov2}
Oquab, M., et al. (2023). DINOv2: Learning robust visual features without supervision. \textit{arXiv:2304.07193}.

\bibitem{johnson2019billion}
Johnson, J., Douze, M., \& Jegou, H. (2019). Billion-scale similarity search with GPUs (FAISS). \textit{IEEE TBD}, 7(3), 535--547.

\bibitem{kaggle_ghim10k}
Kaggle Dataset. (2020). GHIM-10K Image Dataset. \textit{Kaggle}, \url{https://www.kaggle.com/datasets/mcpython/ghim-10k}.

\bibitem{grigorova2007content}
Grigorova, A., et al. (2007). Content-based image retrieval using feature relevance estimation. \textit{Pattern Recognition Letters}, 28(16), 2314--2324.

\end{thebibliography}
